%  =============================================================================
%  Datei-Hinweise & Konfiguration: siehe config.md
%  -> Abschnitt "anhang/ki-prompts/tool-kimi.tex"
%
%  BESONDERHEIT DIESES ABSCHNITTS: Die Prompts wurden in ENGLISCHER Sprache
%  gestellt und sind hier in deutscher Uebersetzung wiedergegeben. Der Hinweis
%  steht sowohl im Vorspann des Abschnitts als auch an jedem einzelnen Prompt
%  ("aus dem Englischen uebersetzt") -- beides beibehalten, damit kein Prompt
%  fuer sich betrachtet als deutschsprachige Eingabe missverstanden wird.
%
%  ABSCHNITTSNUMMER: Dieser Abschnitt ist Anhang III.IV. Das \note-Feld des
%  Bib-Eintrags moonshotKimi2026 verweist darauf. Wird die Reihenfolge der
%  \input-Zeilen in anhang/ki-prompts.tex geaendert, ist der Verweis dort
%  mitzufuehren.
%  =============================================================================

\section{Moonshot AI Kimi (K3)}
\label{sec:prompts-kimi}

\kiwerkzeug{Moonshot AI Kimi, Modell K3}{Modellbezeichner
\texttt{kimi-k3}, Modellvorstellung vom 17.07.2026}{https://www.kimi.com}
% Eintrag im KI-Verzeichnis (Leitfaden Kap. 6.2.4) erzwingen mittels \nocite.
\nocite{moonshotKimi2026}

Eingesetzt für ein allgemeines Codereview der beigelegten Quellcodebeilage
(\cref{cha:anhang-archiv}) sowie für den Erstentwurf des
Änderungsverzeichnisses der Beilage. Das Modell erhielt ausschließlich das
entpackte Archiv und keinen Zugriff auf den Projektordner; es war damit die
einzige Prüfinstanz, die den Code in genau der Form vorfand, in der er einem
Dritten vorliegt.

Die Prompts dieses Abschnitts wurden in englischer Sprache gestellt und sind
nachfolgend sinngleich ins Deutsche übersetzt wiedergegeben; der Zusatz
\textit{aus dem Englischen übersetzt} weist an jedem Prompt darauf hin. Die
Übersetzung war nicht Teil der Eingabe an das Modell.

Wie bei den übrigen Werkzeugen wurde jede Antwort eigenständig gegengeprüft.
Das Ergebnis fiel hier geteilt aus und ist deshalb ausdrücklich vermerkt: Von
zwei vorgeschlagenen Änderungen wurde die erste nach eigener Messung
übernommen -- sie betraf den Selbstschutz des Sitzungserhalts und beschrieb
einen zuvor unentdeckten Fehler --, die zweite als Rückschritt verworfen. Ihr
lag eine Verwechslung zweier Spalten des Archivmanifests zugrunde, die ohne
Kenntnis des Bauprozesses nicht erkennbar war; ihre Übernahme hätte den Bau
der Beilage abgebrochen. Auch der Erstentwurf des Änderungsverzeichnisses
wurde vor der Übernahme in drei Punkten korrigiert.

\subsection{Allgemeines Codereview der Quellcodebeilage}
\label{sec:prompts-kimi-codereview}
\begin{enumerate}
    \item \enquote{Anbei der vollständige Quellcodebestand eines
            Handelsagenten in der Fassung, die einer wissenschaftlichen
            Arbeit als Beilage beigefügt wird. Prüfe ihn ohne Vorgabe eines
            Schwerpunkts auf Fehler und benenne je Befund die betroffene
            Datei, die Zeile, die Wirkung im Betrieb und deine Sicherheit über
            den Befund: [\dots]}
            \par\hfill\textit{Kontext: \cref{cha:anhang-archiv}; aus dem
            Englischen übersetzt}

    \item \enquote{Beschränke die Prüfung nun auf die Shell-Skripte und dort
            auf die Behandlung von Rückgabewerten: Welche Abfrage deutet einen
            Fehlerausgang als regulären Zustand, und in welchem Fall bleibt
            eine notwendige Handlung dadurch unbemerkt aus? [\dots]}
            \par\hfill\textit{Kontext: \cref{cha:anhang-archiv},
            \cref{sec:umsetzung-implementierung}; aus dem Englischen
            übersetzt}

    \item \enquote{Lege zu jedem übernahmewürdigen Befund eine minimale
            Änderung als Textunterschied vor, die die bestehende Struktur
            wahrt, und benenne für jede Änderung, in welchen Umgebungen sie
            das Verhalten ändert und in welchen sie nachweislich wirkungslos
            bleibt: [\dots]}
            \par\hfill\textit{Kontext: \cref{cha:anhang-archiv}; aus dem
            Englischen übersetzt}
\end{enumerate}

\subsection{Entwurf des Änderungsverzeichnisses der Beilage}
\label{sec:prompts-kimi-changelog}
\begin{enumerate}
    \item \enquote{Formuliere zu den bestätigten Befunden einen Abschnitt für
            den Anhang, der die Unterschiede zwischen der vorliegenden und der
            nächsten Version der Beilage benennt: je Punkt das fehlerhafte
            Verhalten davor, die Änderung und den Beleg, dass sie wirkt --
            sachlich, ohne Wertung und ohne Wiedergabe von Quelltext:
            [\dots]}
            \par\hfill\textit{Kontext:
            Anhang~\ref{sec:anhang-archiv-identitaet}; aus dem Englischen
            übersetzt}

    \item \enquote{Prüfe den Abschnitt gegen die Angaben der Beilage: Welche
            Zahl, welcher Dateiname und welcher Verweis darin sind aus dem
            beigefügten Bestand nicht belegbar? [\dots]}
            \par\hfill\textit{Kontext:
            Anhang~\ref{sec:anhang-archiv-identitaet}; aus dem Englischen
            übersetzt}
\end{enumerate}

%  =============================================================================
%  Ende
%  =============================================================================
